%=============================================================================
% AGENT 8: FERRITE-SUPERCONDUCTOR SHELL DESIGN FOR EARTH ATMOSPHERIC OPERATION
% Engineering specification for a psi-bubble device at the Karman line
%=============================================================================

\documentclass[11pt,a4paper]{article}
\usepackage{amsmath,amssymb,amsthm}
\usepackage{booktabs}
\usepackage{graphicx}
\usepackage{siunitx}
\usepackage{geometry}
\usepackage{multirow}
\usepackage{array}
\usepackage{enumitem}
\geometry{margin=2.5cm}

\newtheorem{theorem}{Theorem}
\newtheorem{proposition}{Proposition}
\newtheorem{corollary}{Corollary}
\newtheorem{designrule}{Design Rule}

\title{Ferrite--Superconductor Shell Design\\for Earth Atmospheric $\psi$-Bubble Operation\\
\large Engineering Specification at the K\'arm\'an Line ($h = \SI{100}{km}$)}
\author{Agent 8 -- DFD Research Iteration}
\date{March 2026}

\begin{document}
\maketitle

\begin{abstract}
We present a complete engineering design for a ferrite--superconductor composite
shell optimised for $\psi$-bubble activation at and above the K\'arm\'an line
($h = \SI{100}{km}$), where the ambient mass density
$\rho = \SI{5.6e-7}{kg/m^3}$ requires a threshold magnetic field of only
$B_c = \SI{11.4}{T}$. The design employs REBCO (rare-earth barium copper oxide)
high-temperature superconducting tape wound into a prolate spheroidal shell
generating a persistent \SI{15}{T} field, providing a safety margin of $1.3$
over threshold. We analyse three shell architectures in order of decreasing
complexity: (i) a full ferrite--superconductor composite with YIG resonant
storage for parametric threshold modulation; (ii) a thin-ferrite hybrid for
reduced mass; and (iii) a pure superconductor shell with no ferrite, identified
as the \emph{minimum viable $\psi$-bubble} (MVB). The pure-SC design achieves a
total mass of $\sim\!\SI{37000}{kg}$, generates Mode~I thrust using ambient
thermospheric gas with an exhaust velocity of $\SI{1.8e7}{m/s}$ ($0.06c$) and a
specific impulse of $\SI{1.8e6}{s}$, and requires only $\sim\!\SI{10}{kW}$ of
electrical power for cryogenics and control. We provide detailed mass budgets,
thermal analysis, structural loading calculations, launch-to-activation
trajectories, and a technology readiness assessment for each subsystem.
\end{abstract}

\tableofcontents
\newpage

%=============================================================================
\section{Design Point Selection}\label{sec:design_point}
%=============================================================================

\subsection{The $\eta_c$ Threshold at Altitude}

The DFD threshold condition requires:
\begin{equation}\label{eq:eta_threshold}
\eta \equiv \frac{U_{\mathrm{EM}}}{\rho c^2} \geq \eta_c = \frac{\alpha}{4}
\approx 1.82 \times 10^{-3},
\end{equation}
where $U_{\mathrm{EM}} = B^2/(2\mu_0)$ for a purely magnetic configuration. Solving
for the threshold field:
\begin{equation}\label{eq:B_threshold}
B_{\mathrm{threshold}}(h) = \sqrt{2\mu_0 \,\rho(h)\, c^2 \,\eta_c}\,.
\end{equation}

\subsection{Altitude Trade Study}

Using the US Standard Atmosphere (1976) and NRLMSISE-00 thermospheric model:

\begin{table}[h]
\centering
\caption{Threshold magnetic field as a function of altitude.}
\label{tab:altitude_trade}
\begin{tabular}{@{}lrrrl@{}}
\toprule
Altitude & $\rho$ [\si{kg/m^3}] & $B_{\mathrm{threshold}}$ [\si{T}] & Margin at \SI{15}{T} & Technology \\
\midrule
\SI{0}{km} (sea level) & $1.225$ & $22{,}500$ & $6.7\times 10^{-4}$ & Impossible \\
\SI{40}{km} (balloon ceiling) & $3.99 \times 10^{-3}$ & $1{,}280$ & $0.012$ & Impossible \\
\SI{60}{km} & $3.10 \times 10^{-4}$ & $357$ & $0.042$ & Impossible \\
\SI{80}{km} (mesopause) & $1.8 \times 10^{-5}$ & $86$ & $0.17$ & Pulsed HTS \\
\SI{100}{km} (K\'arm\'an) & $5.6 \times 10^{-7}$ & $11.4$ & $\mathbf{1.3}$ & \textbf{Persistent HTS} \\
\SI{120}{km} & $2.2 \times 10^{-8}$ & $3.0$ & $5.0$ & HTS (easy) \\
\SI{200}{km} (LEO) & $2.5 \times 10^{-10}$ & $0.24$ & $63$ & NdFeB permanent \\
\SI{400}{km} (ISS) & $2.8 \times 10^{-12}$ & $0.025$ & $600$ & Ferrite alone \\
\bottomrule
\end{tabular}
\end{table}

\begin{designrule}[Primary Design Point]
The K\'arm\'an line ($h = \SI{100}{km}$) is selected as the primary operating
altitude because:
\begin{enumerate}[nosep]
\item $B_{\mathrm{threshold}} = \SI{11.4}{T}$ is within the demonstrated
capability of REBCO superconducting magnets (record: \SI{45.5}{T}, NHMFL 2019);
\item a design field of $B_0 = \SI{15}{T}$ provides a safety margin of $1.3$;
\item the altitude is above the homopause, so residual gas is predominantly
atomic oxygen (beneficial for plasma Mode~I thrust);
\item conventional launch vehicles can deliver payloads to \SI{100}{km} routinely.
\end{enumerate}
\end{designrule}

\subsection{Exhaust Velocity and Specific Impulse at the Design Point}

At $B = \SI{15}{T}$ and $\rho = \SI{5.6e-7}{kg/m^3}$, the Alfv\'en speed
(governing Mode~I exhaust velocity) is:
\begin{equation}\label{eq:v_exh}
v_{\mathrm{exh}} = \frac{B}{\sqrt{\mu_0 \rho}}
= \frac{15}{\sqrt{1.257 \times 10^{-6} \times 5.6 \times 10^{-7}}}
= \SI{1.79e7}{m/s} \approx 0.06\,c\,.
\end{equation}
The corresponding specific impulse:
\begin{equation}
I_{\mathrm{sp}} = \frac{v_{\mathrm{exh}}}{g_0}
= \frac{1.79 \times 10^7}{9.81} = \SI{1.82e6}{s}\,.
\end{equation}
For comparison, chemical rockets achieve $I_{\mathrm{sp}} \sim 300$--$450\;\mathrm{s}$,
ion thrusters $\sim 3000$--$10{,}000\;\mathrm{s}$. The $\psi$-bubble Mode~I drive
exceeds all existing propulsion by three orders of magnitude.

%=============================================================================
\section{Shell Architecture: Three Configurations}\label{sec:architectures}
%=============================================================================

We analyse three shell configurations in order of decreasing complexity and
mass, identifying the optimal trade between performance and buildability.

%-----------------------------------------------------------------------------
\subsection{Configuration A: Full Ferrite--Superconductor Composite}
%-----------------------------------------------------------------------------

\subsubsection{Concept}

The full composite shell comprises three functional layers:

\begin{enumerate}[nosep]
\item \textbf{Outer structural shell:} Carbon fibre reinforced polymer (CFRP)
or Inconel 718, providing meteoroid shielding and mechanical support.

\item \textbf{Ferrite resonant layer:} Yttrium iron garnet (YIG,
\ce{Y3Fe5O12}) tiles bonded to the structural shell, providing high-$Q$
($>10^4$) electromagnetic resonant storage for parametric threshold modulation.

\item \textbf{Inner superconductor layer:} REBCO tape (YBCO on Hastelloy
substrate) wound in a cos-$\theta$ or saddle-coil configuration, generating
the persistent \SI{15}{T} base field.

\item \textbf{Innermost radiation/thermal shield:} Multi-layer insulation (MLI)
and cryogenic thermal bus.
\end{enumerate}

\subsubsection{Ferrite Layer Specification}

\begin{table}[h]
\centering
\caption{YIG ferrite layer properties.}
\label{tab:yig_properties}
\begin{tabular}{@{}lr@{}}
\toprule
Property & Value \\
\midrule
Material & \ce{Y3Fe5O12} (YIG) \\
Density & \SI{5170}{kg/m^3} \\
Saturation magnetisation $M_s$ & \SI{0.178}{T} \\
Curie temperature & \SI{553}{K} \\
Unloaded $Q$ at \SI{10}{GHz} & $>10{,}000$ \\
FMR linewidth $\Delta H$ & \SI{0.3}{Oe} (\SI{24}{A/m}) \\
Relative permeability $\mu_r$ & $\sim 15$ (below saturation) \\
Thermal conductivity & \SI{7.4}{W/(m\cdot K)} \\
\bottomrule
\end{tabular}
\end{table}

The ferrite layer serves a specific purpose: storing the \emph{modulation}
energy for parametric pumping of the $\psi$-field. In this mode, the magnetic
field oscillates between $B_{\min} = \SI{13}{T}$ and $B_{\max} = \SI{17}{T}$
at a frequency matched to the $\psi$-field resonance. The YIG tiles convert
between magnetic and electromagnetic energy with minimal loss ($Q > 10^4$),
enabling sustained parametric amplification.

\textbf{Modulation energy per cycle:}
\begin{equation}
\Delta U = \frac{B_{\max}^2 - B_{\min}^2}{2\mu_0} \times V_{\mathrm{cavity}}
= \frac{17^2 - 13^2}{2 \times 1.257 \times 10^{-6}} \times \frac{4}{3}\pi(5)^3
= \SI{2.5e10}{J}\,.
\end{equation}

\subsubsection{Mass Budget: Configuration A}

\begin{table}[h]
\centering
\caption{Mass budget for Configuration A (full ferrite--SC composite).}
\label{tab:mass_A}
\begin{tabular}{@{}lrrr@{}}
\toprule
Subsystem & Unit mass & Quantity/Area & Total mass [\si{kg}] \\
\midrule
REBCO SC coils & \SI{300}{kg/m^3} & \SI{50}{m^3} & $15{,}000$ \\
YIG ferrite tiles (\SI{50}{mm}) & \SI{5170}{kg/m^3} & \SI{17.6}{m^3} & $91{,}000$ \\
CFRP structure & \SI{1600}{kg/m^3} & \SI{9.4}{m^3} & $15{,}000$ \\
Cryogenics (pulse tubes + bus) & -- & -- & $1{,}600$ \\
Power system (solar + batteries) & -- & -- & $500$ \\
Avionics and control & -- & -- & $400$ \\
Payload/crew volume & -- & -- & $5{,}000$ \\
\midrule
\textbf{Total} & & & $\mathbf{128{,}500}$ \\
\bottomrule
\end{tabular}
\end{table}

\begin{designrule}[Ferrite Mass Problem]
At \SI{91}{t}, the YIG ferrite layer constitutes 71\% of the total vehicle
mass. This is the dominant engineering challenge for Configuration~A.
\end{designrule}

%-----------------------------------------------------------------------------
\subsection{Configuration B: Thin-Ferrite Hybrid}
%-----------------------------------------------------------------------------

\subsubsection{Rationale for Reduced Ferrite}

The full \SI{50}{mm} ferrite thickness was sized for maximum modulation
depth. However, the parametric pumping can operate with reduced modulation
amplitude at the cost of longer $\psi$-field growth time. A \SI{5}{mm}
ferrite shell reduces the modulation energy by $10\times$ but the $Q$-factor
remains $>10^4$, so the exponential growth rate decreases by only $\ln(10)
\approx 2.3\times$.

\subsubsection{Mass Budget: Configuration B}

\begin{table}[h]
\centering
\caption{Mass budget for Configuration B (thin-ferrite hybrid).}
\label{tab:mass_B}
\begin{tabular}{@{}lr@{}}
\toprule
Subsystem & Total mass [\si{kg}] \\
\midrule
REBCO SC coils & $15{,}000$ \\
YIG ferrite tiles (\SI{5}{mm}) & $9{,}100$ \\
CFRP structure & $15{,}000$ \\
Cryogenics & $1{,}600$ \\
Power system & $500$ \\
Avionics and control & $400$ \\
Payload & $5{,}000$ \\
\midrule
\textbf{Total} & $\mathbf{46{,}600}$ \\
\bottomrule
\end{tabular}
\end{table}

Configuration~B reduces total mass by 64\% relative to Configuration~A while
retaining parametric modulation capability. The ferrite fraction drops from
71\% to 20\%.

%-----------------------------------------------------------------------------
\subsection{Configuration C: Pure Superconductor Shell (Minimum Viable $\psi$-Bubble)}
%-----------------------------------------------------------------------------

\subsubsection{Rationale}

For \emph{static} threshold operation -- where $B_0 > B_c$ is maintained as a
persistent superconducting current with no parametric modulation -- the ferrite
layer is not required. The $\psi$-field transition is triggered by simply
exceeding $\eta_c$; no resonant pumping is needed. This eliminates the heaviest
subsystem entirely.

\begin{designrule}[Minimum Viable $\psi$-Bubble]
If the DFD threshold transition occurs spontaneously when $\eta > \eta_c$
(as predicted by the theory), then no parametric resonance is required and
the ferrite layer can be eliminated. Configuration~C is the minimum viable
$\psi$-bubble (MVB).
\end{designrule}

\subsubsection{Mass Budget: Configuration C}

\begin{table}[h]
\centering
\caption{Mass budget for Configuration C (pure SC shell, MVB).}
\label{tab:mass_C}
\begin{tabular}{@{}lrl@{}}
\toprule
Subsystem & Mass [\si{kg}] & Notes \\
\midrule
REBCO SC coils & $15{,}000$ & 50 layers, \SI{5}{mm} total \\
CFRP structural shell & $12{,}000$ & Reduced without ferrite load \\
Cryogenics & $1{,}600$ & 4 pulse-tube cryocoolers \\
Power system & $500$ & Solar array + Li-ion buffer \\
Avionics/GNC & $400$ & Redundant flight computers \\
Thermal management & $1{,}500$ & MLI + radiators \\
Crew/payload & $5{,}000$ & Pressurised cabin volume \\
Launch adapter & $1{,}000$ & Jettisoned after separation \\
\midrule
\textbf{Total (on-orbit)} & $\mathbf{37{,}000}$ & After adapter jettison \\
\bottomrule
\end{tabular}
\end{table}

%=============================================================================
\section{Superconductor Subsystem Design}\label{sec:sc_design}
%=============================================================================

\subsection{REBCO Tape Selection}

\begin{table}[h]
\centering
\caption{REBCO tape baseline specifications (representative of SuperPower
2G-HTS or Fujikura FESC).}
\label{tab:rebco_spec}
\begin{tabular}{@{}lr@{}}
\toprule
Parameter & Value \\
\midrule
Conductor architecture & YBCO on Hastelloy C-276 \\
Tape width & \SI{12}{mm} \\
Tape thickness & \SI{0.1}{mm} (including substrate) \\
Critical current $I_c$ (\SI{77}{K}, self-field) & \SI{500}{A} \\
Critical current $I_c$ (\SI{20}{K}, \SI{15}{T}) & $\sim\!\SI{250}{A}$ \\
Engineering current density $J_e$ (\SI{20}{K}, \SI{15}{T}) & $\sim\!\SI{200}{A/mm^2}$ \\
Irreversibility field $B^*$ (\SI{20}{K}) & $>\SI{40}{T}$ \\
Operating temperature range & \SIrange{4}{77}{K} \\
Minimum bend radius & \SI{5.5}{mm} \\
Mechanical stress limit & \SI{700}{MPa} \\
Piece length (commercial) & $>\SI{300}{m}$ \\
\bottomrule
\end{tabular}
\end{table}

\subsection{Coil Geometry: Prolate Spheroidal Shell}

The shell is a prolate spheroid with semi-major axis $a = \SI{6}{m}$ and
semi-minor axis $b = \SI{5}{m}$ (eccentricity $e = 0.553$). The REBCO tape
is wound in a modified cos-$\theta$ pattern to produce a predominantly dipolar
field inside the cavity, with field concentration at the poles for directional
thrust vectoring.

\textbf{Shell surface area:}
\begin{equation}
A = 2\pi b^2 + \frac{2\pi ab}{e}\sin^{-1}(e)
= 2\pi(25) + \frac{2\pi(30)}{0.553}\sin^{-1}(0.553)
= \SI{352}{m^2}\,.
\end{equation}

\textbf{REBCO tape volume} (50 layers $\times$ \SI{0.1}{mm} thickness):
\begin{equation}
V_{\mathrm{SC}} = A \times 50 \times 0.1\;\mathrm{mm} = 352 \times 0.005 = \SI{1.76}{m^3}\,.
\end{equation}

\textbf{REBCO tape mass} (density of Hastelloy-dominated tape $\approx \SI{8500}{kg/m^3}$):
\begin{equation}
m_{\mathrm{SC}} \approx 1.76 \times 8500 \approx \SI{15000}{kg}\,.
\end{equation}

\textbf{Total tape length:}
\begin{equation}
L_{\mathrm{tape}} = \frac{A \times 50}{0.012} = \frac{352 \times 50}{0.012}
= \SI{1.47e6}{m} = \SI{1470}{km}\,.
\end{equation}
This is a large but not unprecedented quantity. The ITER tokamak uses
$\sim\!\SI{100000}{km}$ of Nb$_3$Sn and NbTi conductor.

\subsection{Persistent Current Mode}

Once energised, the REBCO coil operates in persistent current mode with
zero external power. The persistent current switch is a short section of
the REBCO tape heated above $T_c$ during charging, then cooled to close
the circuit.

\textbf{Stored magnetic energy:}
\begin{equation}
E_{\mathrm{mag}} = \frac{B^2}{2\mu_0} V_{\mathrm{cavity}}
\approx \frac{15^2}{2 \times 1.257 \times 10^{-6}} \times \frac{4}{3}\pi(4.5)^3
\approx \SI{3.4e10}{J} = \SI{34}{GJ}\,.
\end{equation}
This is equivalent to about 8 tonnes of TNT. The energy is stored entirely
in the magnetic field and can be recovered by de-energising the coil.

\textbf{Charging time} (at \SI{10}{MW} power supply, ground-based):
\begin{equation}
t_{\mathrm{charge}} = \frac{E_{\mathrm{mag}}}{P_{\mathrm{supply}}}
= \frac{3.4 \times 10^{10}}{10^7} = \SI{3400}{s} \approx \SI{57}{min}\,.
\end{equation}

\subsection{Meissner Shielding of the Interior}

The superconducting shell provides complete Meissner shielding of the interior
cavity. This creates:
\begin{itemize}[nosep]
\item A field-free region for crew and electronics (important for radiation-
sensitive components);
\item A sharp field boundary at the shell surface, maximising the field
gradient and thus the $\psi$-field coupling;
\item Natural containment of the $\psi$-bubble within the shell geometry.
\end{itemize}

The London penetration depth for YBCO is $\lambda_L \approx \SI{150}{nm}$,
so the field decays to $<10^{-6} B_0$ within a few micrometres of the inner
surface. With 50 layers of tape, the shielding is effectively perfect.

%=============================================================================
\section{Shell Geometry and Structural Analysis}\label{sec:structural}
%=============================================================================

\subsection{Magnetic Pressure Loading}

The magnetic pressure on the shell inner surface is:
\begin{equation}
P_{\mathrm{mag}} = \frac{B^2}{2\mu_0} = \frac{15^2}{2 \times 1.257 \times 10^{-6}}
= \SI{8.95e7}{Pa} \approx \SI{90}{MPa}\,.
\end{equation}

This is a substantial outward pressure (equivalent to $\sim\!900\;\mathrm{atm}$)
that the structural shell must resist.

\subsection{Hoop Stress in the Spheroidal Shell}

For a thin-walled pressure vessel with semi-minor axis $b = \SI{5}{m}$ and
wall thickness $t = \SI{0.1}{m}$ (combined SC + structure):
\begin{equation}
\sigma_{\mathrm{hoop}} = \frac{P_{\mathrm{mag}} \cdot b}{t}
= \frac{8.95 \times 10^7 \times 5}{0.1} = \SI{4.48}{GPa}\,.
\end{equation}

\begin{table}[h]
\centering
\caption{Structural material options for the pressure shell.}
\label{tab:structural_materials}
\begin{tabular}{@{}lrrl@{}}
\toprule
Material & UTS [\si{GPa}] & Density [\si{kg/m^3}] & Feasibility \\
\midrule
CFRP (T1000G/epoxy) & 3.0 & 1600 & Insufficient alone \\
Inconel 718 & 1.4 & 8200 & Insufficient alone \\
Toray T1100G CFRP & 3.7 & 1790 & Marginal \\
\textbf{REBCO tape (Hastelloy)} & \textbf{1.2} & \textbf{8500} & \textbf{Self-supporting to \SI{700}{MPa}} \\
\bottomrule
\end{tabular}
\end{table}

\begin{designrule}[Structural Solution]
The \SI{4.5}{GPa} hoop stress exceeds any single material's UTS at reasonable
wall thickness. The solution is to increase wall thickness or employ external
reinforcement:
\begin{enumerate}[nosep]
\item Increase structural shell to $t = \SI{0.3}{m}$: reduces $\sigma$ to
\SI{1.5}{GPa} (achievable with high-modulus CFRP layup);
\item Use the REBCO tape itself as structural reinforcement: the Hastelloy
substrate carries \SI{700}{MPa}, contributing to the overall load path;
\item Employ \emph{force-balanced} coil design where opposing field regions
cancel the net outward force (as in tokamak designs).
\end{enumerate}
\end{designrule}

The force-balanced approach is the most mass-efficient. In a dipole
configuration, the attractive force between opposing-current paths partially
offsets the expansive magnetic pressure. The net structural load is reduced
by approximately 40\%, bringing the required wall thickness to $\sim\!\SI{0.2}{m}$
and the hoop stress to $\sim\!\SI{2.2}{GPa}$ -- achievable with advanced CFRP.

%=============================================================================
\section{Thermal Management and Cryogenics}\label{sec:cryo}
%=============================================================================

\subsection{Operating Temperature}

REBCO tape at \SI{15}{T} has a comfortable operating margin at \SI{20}{K}:
\begin{itemize}[nosep]
\item $T_c$ at \SI{15}{T}: $\sim\!\SI{75}{K}$
\item Operating temperature: \SI{20}{K}
\item Temperature margin: \SI{55}{K} (exceptionally large)
\item $I_c$ at \SI{20}{K}, \SI{15}{T}: $\sim\!\SI{250}{A}$ per \SI{12}{mm} tape
\end{itemize}

\subsection{Heat Load Budget}

\begin{table}[h]
\centering
\caption{Cryogenic heat load at \SI{20}{K}.}
\label{tab:heat_load}
\begin{tabular}{@{}lrl@{}}
\toprule
Source & Heat load [\si{W}] & Notes \\
\midrule
Radiation through MLI & 15 & 40 layers MLI, $\epsilon_{\mathrm{eff}} = 10^{-4}$ \\
Conduction through supports & 8 & G-10 fibreglass standoffs \\
Current lead heat leak & 0 & Persistent mode (leads retracted) \\
Residual AC loss (vibration) & 2 & Magnetostrictive coupling \\
Instrumentation wiring & 1 & Phosphor-bronze wires \\
\midrule
\textbf{Total at \SI{20}{K}} & \textbf{26} & \\
\bottomrule
\end{tabular}
\end{table}

\subsection{Cryocooler Selection}

Pulse-tube cryocoolers are selected for their vibration-free operation
(no moving parts at the cold end):

\begin{itemize}[nosep]
\item Model class: Cryomech PT420 or equivalent
\item Cooling capacity: $\sim\!\SI{14}{W}$ at \SI{20}{K} per unit
\item Number of units: 4 (redundant: 2 required, 4 installed)
\item Electrical power per unit: \SI{10}{kW} compressor input
\item Total cryocooler mass: \SI{400}{kg} per unit $\times 4 = \SI{1600}{kg}$
\item Total electrical power for cryo: \SI{20}{kW} (2 running) to \SI{40}{kW} (all 4)
\end{itemize}

\textbf{Note:} In the vacuum of \SI{100}{km} altitude, the radiative heat
load is dominated by solar input on the sunlit side and thermal emission
on the dark side. Active thermal control (sun shields + radiators) reduces
the effective load to the values in Table~\ref{tab:heat_load}.

%=============================================================================
\section{Power System}\label{sec:power}
%=============================================================================

\subsection{Steady-State Power Requirements}

\begin{table}[h]
\centering
\caption{Electrical power budget.}
\label{tab:power}
\begin{tabular}{@{}lr@{}}
\toprule
Subsystem & Power [\si{kW}] \\
\midrule
Cryocoolers (2 of 4 running) & 5.0 \\
Thrust vectoring RF amplifiers & 2.0 \\
Avionics and GNC & 0.5 \\
Thermal management & 1.0 \\
Life support (if crewed) & 1.5 \\
\midrule
\textbf{Total} & \textbf{10.0} \\
\bottomrule
\end{tabular}
\end{table}

\subsection{Power Source Options}

\begin{enumerate}[nosep]
\item \textbf{Solar arrays:} At \SI{100}{km}+, solar flux is the full
$\SI{1361}{W/m^2}$. With 30\% efficient GaAs triple-junction cells,
\SI{10}{kW} requires $\sim\!\SI{25}{m^2}$ of array area.
Specific power: $\sim\!\SI{200}{W/kg}$ $\Rightarrow$ array mass $\sim\!\SI{50}{kg}$.

\item \textbf{RTG (radioisotope):} For missions where solar is impractical
(e.g., extended interplanetary). A \SI{10}{kW} Stirling RTG would mass
$\sim\!\SI{500}{kg}$.

\item \textbf{Lithium-ion battery buffer:} \SI{100}{kWh} buffer for
peak loads and eclipse periods. At \SI{250}{Wh/kg}, mass $= \SI{400}{kg}$.
\end{enumerate}

\textbf{Selected baseline:} Solar + Li-ion buffer, total mass \SI{500}{kg}.

\subsection{Energisation Power (Ground-Based)}

The \SI{34}{GJ} coil charging is performed on the ground before launch
via an external \SI{10}{MW} power supply (equivalent to a small substation).
This is a one-time operation taking $\sim\!\SI{1}{hour}$. The persistent
current then sustains the field indefinitely with zero on-board power.

%=============================================================================
\section{Thrust Vectoring and Propulsion}\label{sec:thrust}
%=============================================================================

\subsection{Mode I Thrust Mechanism}

At \SI{100}{km}, the ambient thermospheric gas (predominantly atomic oxygen,
$n \approx \SI{5e15}{m^{-3}}$) is swept into the $\psi$-bubble's magnetic
field structure. The DFD coupling accelerates this gas to the Alfv\'en speed
(Eq.~\ref{eq:v_exh}), ejecting it as a collimated exhaust.

\textbf{Mass flow rate} (ambient ram intake at orbital velocity
$v_{\mathrm{orb}} = \SI{7800}{m/s}$, effective capture area $A_c = \SI{100}{m^2}$):
\begin{equation}
\dot{m} = \rho \, v_{\mathrm{orb}} \, A_c
= 5.6 \times 10^{-7} \times 7800 \times 100
= \SI{4.4e-1}{kg/s}\,.
\end{equation}

\textbf{Thrust:}
\begin{equation}
F = \dot{m} \, v_{\mathrm{exh}}
= 0.44 \times 1.79 \times 10^7
= \SI{7.9e6}{N} \approx \SI{8}{MN}\,.
\end{equation}

\textbf{Thrust-to-weight ratio} (at vehicle mass \SI{37000}{kg}):
\begin{equation}
\mathrm{TWR} = \frac{F}{mg} = \frac{7.9 \times 10^6}{37000 \times 9.81} = 21.8\,.
\end{equation}

This is extraordinary. For comparison, the Space Shuttle main engines
produced a TWR of $\sim\!1.5$.

\subsection{Thrust Vectoring}

Directional thrust is achieved by breaking the symmetry of the magnetic
field distribution:
\begin{itemize}[nosep]
\item \textbf{Asymmetric coil currents:} Increasing current in one hemisphere
concentrates the field (and thus the exhaust) in that direction;
\item \textbf{Ferrite modulation} (Configurations A and B): Selectively
energising ferrite sectors rotates the field asymmetry;
\item \textbf{External steering coils:} Small superconducting trim coils
on the shell surface provide fine pointing control.
\end{itemize}

For Configuration~C (no ferrite), thrust vectoring relies on asymmetric
coil currents and trim coils. Response time is limited by the $L/R$ time
constant of the SC circuit ($\sim\!\SI{1}{s}$ with resistive heater switches).

%=============================================================================
\section{Launch-to-Activation Trajectory}\label{sec:launch}
%=============================================================================

The $\psi$-bubble device must be transported to \SI{100}{km} altitude before
the magnetic field can cross threshold. Four options:

\subsection{Option A: Balloon + Rocket}

\begin{enumerate}[nosep]
\item High-altitude balloon lifts vehicle to \SI{40}{km} (requires
\SI{37000}{kg} lift capacity -- possible with cluster balloon array);
\item Solid or hybrid rocket stage accelerates from \SI{40}{km} to
\SI{100}{km} (requires $\Delta v \approx \SI{1400}{m/s}$);
\item At \SI{100}{km}, $\psi$-bubble activates; Mode~I thrust takes over.
\end{enumerate}
\textbf{Pros:} Lowest propellant mass. \textbf{Cons:} Cluster balloon
at \SI{37}{t} is at the edge of feasibility.

\subsection{Option B: Electromagnetic Launch}

\begin{enumerate}[nosep]
\item Ground-based electromagnetic launcher (railgun or coilgun) accelerates
vehicle at $\sim\!100\,g$ over \SI{2}{km} track to $v = \SI{2000}{m/s}$;
\item Ballistic coast to \SI{100}{km};
\item $\psi$-bubble activates at apogee.
\end{enumerate}
\textbf{Pros:} No propellant. \textbf{Cons:} Extreme $g$-loads prohibit
crew; requires multi-GW pulsed power supply.

\subsection{Option C: Conventional Rocket (Baseline)}

\begin{enumerate}[nosep]
\item Standard expendable or reusable launch vehicle (e.g., Falcon~9 class)
delivers vehicle to \SI{100}{km} $\times$ \SI{200}{km} suborbital trajectory;
\item At \SI{100}{km}, persistent SC field already active ($B_0 = \SI{15}{T}$,
pre-charged on ground);
\item $\eta > \eta_c$ condition met spontaneously; $\psi$-bubble forms;
\item Mode~I thrust provides $\Delta v$ for orbit insertion and beyond.
\end{enumerate}
\textbf{Pros:} Uses existing launch infrastructure; crew-rated possible.
\textbf{Cons:} Highest cost per launch.

\subsection{Option D: Plasma Sheath Bootstrap (Speculative)}

If the plasma sheath bootstrap mechanism described in Agent~2's analysis
functions as predicted, the device could activate at lower altitudes
($h \sim \SI{60}$--$\SI{80}{km}$), relaxing launch requirements. This
option is contingent on validating the self-sustaining plasma density
reduction feedback loop.

\textbf{Baseline selection: Option~C} (conventional rocket to \SI{100}{km}).

%=============================================================================
\section{Above-Threshold Field Envelope}\label{sec:field_envelope}
%=============================================================================

\subsection{Dipole Field Falloff}

The shell generates an approximately dipolar field outside its surface.
The field magnitude at distance $r$ from the centre (along the equatorial
plane) is:
\begin{equation}
B(r) = B_0 \left(\frac{R}{r}\right)^3, \quad r \geq R,
\end{equation}
where $R = \SI{5}{m}$ (semi-minor axis, equatorial surface).

\subsection{Extent of the Above-Threshold Region}

The $\psi$-bubble forms where $\eta > \eta_c$. At \SI{100}{km}, the threshold
field is $B_c = \SI{11.4}{T}$. The above-threshold region extends to:
\begin{equation}
r_{\mathrm{max}} = R \left(\frac{B_0}{B_c}\right)^{1/3}
= 5 \times \left(\frac{15}{11.4}\right)^{1/3}
= 5 \times 1.095 = \SI{5.5}{m}\,.
\end{equation}

With only a $1.3\times$ margin, the above-threshold region extends only
\SI{0.5}{m} beyond the shell surface. This is adequate for $\psi$-bubble
nucleation but not for a large-volume bubble.

\textbf{To extend the bubble:} Increase $B_0$ or operate at higher altitude.
At \SI{200}{km} ($B_c = \SI{0.24}{T}$):
\begin{equation}
r_{\mathrm{max}} = 5 \times \left(\frac{15}{0.24}\right)^{1/3}
= 5 \times 3.97 = \SI{19.8}{m}\,.
\end{equation}

At LEO altitudes, the $\psi$-bubble extends $\sim\!4R$ from the shell,
providing a substantial interaction volume for Mode~I thrust.

%=============================================================================
\section{Technology Readiness Assessment}\label{sec:tra}
%=============================================================================

\begin{table}[h]
\centering
\caption{Technology readiness levels (TRL) for shell subsystems.}
\label{tab:trl}
\begin{tabular}{@{}lrl@{}}
\toprule
Subsystem & TRL & Justification \\
\midrule
REBCO tape manufacturing & 7--8 & Commercial production; >100 km/yr \\
\SI{15}{T} SC magnets & 6--7 & NHMFL achieved \SI{45.5}{T}; MRI at \SI{11.7}{T} \\
Persistent current switches & 5--6 & Demonstrated in lab SC magnets \\
Prolate spheroid winding & 3--4 & Novel geometry; flat-race winding at TRL 5 \\
Pulse-tube cryocoolers & 8--9 & Space-qualified (JWST heritage) \\
CFRP pressure vessel (\SI{2}{GPa}) & 4--5 & Demonstrated at smaller scale \\
YIG ferrite tiles (large area) & 5--6 & Used in radar/telecom; not at this scale \\
$\psi$-bubble physics & 1--2 & \textbf{Theoretical prediction, unverified} \\
Mode~I thrust mechanism & 1--2 & \textbf{Theoretical prediction, unverified} \\
\midrule
\textbf{Overall system} & \textbf{2} & Limited by physics validation \\
\bottomrule
\end{tabular}
\end{table}

\begin{designrule}[Critical Path]
Every engineering subsystem is at TRL~3 or above. The system-level TRL
is limited entirely by the unverified DFD physics (TRL~1--2). The critical
path to a flyable $\psi$-bubble device is therefore:
\begin{enumerate}[nosep]
\item Validate the $\eta_c$ threshold in laboratory experiment (TRL 1
$\to$ 3);
\item Demonstrate $\psi$-field formation in a \SI{15}{T} SC magnet
with controlled low-density gas (TRL 3 $\to$ 5);
\item Build subscale (\SI{0.5}{m}) shell prototype and measure thrust
(TRL 5 $\to$ 6);
\item Full-scale ground test with vacuum chamber (TRL 6 $\to$ 7);
\item Suborbital flight test (TRL 7 $\to$ 8).
\end{enumerate}
\end{designrule}

%=============================================================================
\section{Configuration Comparison and Recommendation}\label{sec:comparison}
%=============================================================================

\begin{table}[h]
\centering
\caption{Comparison of shell configurations.}
\label{tab:comparison}
\begin{tabular}{@{}lrrr@{}}
\toprule
Parameter & Config.\ A & Config.\ B & Config.\ C \\
& (Full ferrite) & (Thin ferrite) & (SC only) \\
\midrule
Total mass [\si{kg}] & $128{,}500$ & $46{,}600$ & $37{,}000$ \\
Ferrite mass fraction & 71\% & 20\% & 0\% \\
Parametric modulation & Full & Reduced & None \\
$\psi$-field growth mode & Resonant & Resonant (slower) & Spontaneous \\
Thrust vectoring & Ferrite + coils & Ferrite + coils & Coils only \\
Structural complexity & High & Moderate & \textbf{Low} \\
Launch vehicle class & SLS/Starship & Falcon Heavy & \textbf{Falcon Heavy} \\
TRL (engineering) & 3 & 3--4 & \textbf{4} \\
Risk & Moderate & Low--Moderate & \textbf{Lowest} \\
\bottomrule
\end{tabular}
\end{table}

\textbf{Recommendation:} Configuration~C (pure SC shell) is the preferred
first-generation design. It eliminates the dominant mass driver (ferrite),
uses the most mature technology, fits within existing heavy-lift launch
vehicle capabilities, and tests the fundamental $\psi$-bubble physics
without the additional complication of parametric resonance. If
spontaneous threshold crossing proves insufficient and parametric pumping
is required, Configuration~B can be developed as a second-generation
upgrade.

%=============================================================================
\section{Summary of Design Parameters}\label{sec:summary}
%=============================================================================

\begin{table}[h]
\centering
\caption{Configuration C (MVB) design summary.}
\label{tab:final_summary}
\begin{tabular}{@{}ll@{}}
\toprule
\textbf{Parameter} & \textbf{Value} \\
\midrule
\multicolumn{2}{@{}l}{\textit{Geometry}} \\
Shell shape & Prolate spheroid \\
Semi-major axis $a$ & \SI{6}{m} \\
Semi-minor axis $b$ & \SI{5}{m} \\
Shell surface area & \SI{352}{m^2} \\
Inner cavity volume & \SI{524}{m^3} \\
Wall thickness (total) & \SI{200}{mm} \\
\midrule
\multicolumn{2}{@{}l}{\textit{Magnetic field}} \\
Design field $B_0$ & \SI{15}{T} \\
Threshold field $B_c$ at \SI{100}{km} & \SI{11.4}{T} \\
Safety margin & 1.3 \\
Stored energy & \SI{34}{GJ} \\
Field topology & Dipole (cos-$\theta$ winding) \\
\midrule
\multicolumn{2}{@{}l}{\textit{Superconductor}} \\
Material & REBCO (YBCO/Hastelloy) \\
Operating temperature & \SI{20}{K} \\
Number of tape layers & 50 \\
Total tape length & \SI{1470}{km} \\
Current mode & Persistent (zero power) \\
\midrule
\multicolumn{2}{@{}l}{\textit{Propulsion (Mode I at \SI{100}{km})}} \\
Exhaust velocity $v_{\mathrm{exh}}$ & \SI{1.79e7}{m/s} ($0.06c$) \\
Specific impulse $I_{\mathrm{sp}}$ & \SI{1.82e6}{s} \\
Thrust (orbital ram intake) & \SI{8}{MN} \\
Thrust-to-weight ratio & 21.8 \\
\midrule
\multicolumn{2}{@{}l}{\textit{Mass and power}} \\
Total vehicle mass & \SI{37000}{kg} \\
Payload capacity & \SI{5000}{kg} \\
On-board power & \SI{10}{kW} (solar + battery) \\
Ground charging power & \SI{10}{MW} for \SI{1}{hour} \\
\midrule
\multicolumn{2}{@{}l}{\textit{Operations}} \\
Design operating altitude & $\geq \SI{100}{km}$ \\
Launch method & Conventional rocket (Falcon Heavy class) \\
Activation & Spontaneous at $\eta > \eta_c$ \\
\bottomrule
\end{tabular}
\end{table}

%=============================================================================
\section{Conclusions and Forward Work}\label{sec:conclusions}
%=============================================================================

We have presented a detailed engineering design for a ferrite--superconductor
shell optimised for $\psi$-bubble operation in Earth's atmosphere. The key
findings are:

\begin{enumerate}
\item The K\'arm\'an line (\SI{100}{km}) is the optimal activation altitude,
requiring only $B = \SI{11.4}{T}$ -- well within REBCO superconductor
capability.

\item Three shell configurations were analysed. The pure superconductor
shell (Configuration~C) is recommended as the first-generation design at
\SI{37}{t} total mass, eliminating the \SI{91}{t} ferrite mass penalty.

\item The persistent \SI{15}{T} field requires no on-board power to
maintain. Total vehicle power consumption is only \SI{10}{kW}, dominated
by cryocooling.

\item Mode~I thrust at \SI{100}{km} produces \SI{8}{MN} with $I_{\mathrm{sp}}
= \SI{1.8e6}{s}$, exceeding all existing propulsion technologies by
orders of magnitude.

\item The \SI{90}{MPa} magnetic pressure on the shell requires advanced
CFRP structural design at \SI{2}{GPa} hoop stress, which is achievable
with force-balanced coil geometry.

\item All engineering subsystems are at TRL~3+. The system-level TRL is
limited to TRL~2 solely by the unverified DFD threshold physics.

\item The critical path is physics validation, not engineering development.
A laboratory demonstration of $\eta_c$ threshold crossing would immediately
advance the system TRL to 3+.
\end{enumerate}

\textbf{Forward work:}
\begin{itemize}[nosep]
\item Design a tabletop $\eta_c$ threshold experiment using a \SI{20}{T}
Bitter magnet and low-pressure gas cell (see Agent~10 specification);
\item Develop detailed finite-element structural model of the force-balanced
prolate spheroid under \SI{15}{T} loading;
\item Prototype cos-$\theta$ REBCO winding on curved mandrel at subscale;
\item Model the Mode~I gas-dynamic interaction in the thermosphere using
MHD simulation codes;
\item Analyse orbital mechanics for the $\psi$-bubble vehicle: drag, thrust
profiles, orbit-raising trajectories.
\end{itemize}

\end{document}
